\chapter{Kernel Support Vector Machine}
\lecture{19}{4 \& 11 Nov.}{}

\begin{prev}
    The dual SVM (\autoref{sec:dual-messages}) is another QP, with valuable
    geometric messages and \emph{almost} no dependence on \(\tilde{d}\): it
    has \(N\) variables and \(N + 1\) constraints, but every entry \(q_{n,m} =
    y_n y_m \bm{z}_n^\top \bm{z}_m\) of its matrix is still an inner product
    in \(\mathbb{R}^{\tilde{d}}\).
\end{prev}

This chapter removes the last trace of \(\tilde{d}\) by never forming
\(\bm{z}_n\) at all: for the right transforms, the inner product
\(\bm{\Phi}(\bm{x})^\top \bm{\Phi}(\bm{x}')\) can be computed directly from
\(\bm{x}\) and \(\bm{x}'\), much faster than transforming first. That
shortcut, the kernel, makes polynomial transforms of any degree cheap and
even admits a transform of infinite dimension. The chapter ends by comparing
the kernels and asking which functions can be kernels at all.

\paragraph{Roadmap.}
\begin{enumerate}
    \item \textbf{Embedding Numerous Features: Kernel Models}
        \begin{itemize}
            \item[] Lecture 2: Dual Support Vector Machine --- dual SVM:
                another QP with valuable geometric messages and almost no
                dependence on \(\tilde{d}\)
            \item[] \textbf{Lecture 3: Kernel Support Vector Machine} ---
                Kernel Trick; Polynomial Kernel; Gaussian Kernel; Comparison
                of Kernels
        \end{itemize}
    \item Combining Predictive Features: Aggregation Models
    \item Distilling Implicit Features: Extraction Models
\end{enumerate}

\section{Kernel Trick}\label{sec:kernel-trick}

\source{slides p.2--6}

\subsection{Dual SVM Revisited}

Goal: SVM without dependence on \(\tilde{d}\). Half-way done:
\[
    \begin{aligned}
        \min_{\bm{\alpha}} \quad & \frac{1}{2} \bm{\alpha}^\top
            \mathrm{Q}_{\mathrm{D}} \bm{\alpha} - \bm{1}^\top \bm{\alpha}\\
        \text{subject to} \quad & \bm{y}^\top \bm{\alpha} = 0;\\
        & \alpha_n \ge 0, \text{ for } n = 1, 2, \dots , N
    \end{aligned}
\]
\begin{itemize}
    \item \(q_{n,m} = y_n y_m \bm{z}_n^\top \bm{z}_m\): an inner product in
        \(\mathbb{R}^{\tilde{d}}\);
    \item need: \(\bm{z}_n^\top \bm{z}_m = \bm{\Phi}(\bm{x}_n)^\top
        \bm{\Phi}(\bm{x}_m)\) calculated \textbf{faster} than
        \(O(\tilde{d})\).
\end{itemize}

\begin{question}
    Can we do so?
\end{question}

\subsection{Fast Inner Product for \texorpdfstring{\(\bm{\Phi}_2\)}{Phi2}}

Take the 2nd order polynomial transform
\[
    \bm{\Phi}_2(\bm{x}) = (1, x_1, x_2, \dots , x_d, x_1^2, x_1 x_2, \dots ,
    x_1 x_d, x_2 x_1, x_2^2, \dots , x_2 x_d, \dots , x_d^2)
\]
--- which includes both \(x_1 x_2\) \& \(x_2 x_1\) for `simplicity' :-).
Its inner product splits by degree:
\[
    \begin{aligned}
        \bm{\Phi}_2(\bm{x})^\top \bm{\Phi}_2(\bm{x}') &= 1 + \sum_{i=1}^d x_i
            x'_i + \sum_{i=1}^d \sum_{j=1}^d x_i x_j x'_i x'_j\\
        &= 1 + \sum_{i=1}^d x_i x'_i + \sum_{i=1}^d x_i x'_i \sum_{j=1}^d x_j
            x'_j\\
        &= 1 + \bm{x}^\top \bm{x}' + (\bm{x}^\top \bm{x}')(\bm{x}^\top
            \bm{x}').
    \end{aligned}
\]
The duplicated cross terms are what make the double sum factor: with every
ordered pair \((i, j)\) present, \(\sum_{i,j} (x_i x'_i)(x_j x'_j)\) is a
product of two copies of the same sum. Transforming first costs \(O(d^2)\)
for the \(1 + d + d^2\) coordinates; the last line needs one inner product
in \(\mathbb{R}^d\), \(O(d)\).

\begin{moral}
    For \(\bm{\Phi}_2\), transform \(+\) inner product can be carefully done
    in \(O(d)\) instead of \(O(d^2)\).
\end{moral}

\subsection{Kernel: Transform + Inner Product}

\begin{definition}[Kernel function]
    A transform \(\bm{\Phi}\) comes with its \vocab{kernel function}
    \[
        K_{\bm{\Phi}}(\bm{x}, \bm{x}') \equiv \bm{\Phi}(\bm{x})^\top
        \bm{\Phi}(\bm{x}').
    \]
    For example, \(\bm{\Phi}_2 \iff K_{\bm{\Phi}_2}(\bm{x}, \bm{x}') = 1 +
    (\bm{x}^\top \bm{x}') + (\bm{x}^\top \bm{x}')^2\).
\end{definition}

The dual SVM touches \(\bm{z}\) in three places, and each of them is an
inner product of two transformed vectors.
\begin{itemize}
    \item Quadratic coefficient: \(q_{n,m} = y_n y_m \bm{z}_n^\top \bm{z}_m
        = y_n y_m K(\bm{x}_n, \bm{x}_m)\).
    \item Optimal bias \(b\)? From an SV \((\bm{x}_s, y_s)\),
        \[
            b = y_s - \bm{w}^\top \bm{z}_s = y_s - \left( \sum_{n=1}^N
            \alpha_n y_n \bm{z}_n \right)^\top \bm{z}_s = y_s - \sum_{n=1}^N
            \alpha_n y_n K(\bm{x}_n, \bm{x}_s).
        \]
    \item Optimal hypothesis \(g_{\mathrm{SVM}}\): for a test input
        \(\bm{x}\),
        \[
            g_{\mathrm{SVM}}(\bm{x}) = \sign \left( \bm{w}^\top
            \bm{\Phi}(\bm{x}) + b \right) = \sign \left( \sum_{n=1}^N
            \alpha_n y_n K(\bm{x}_n, \bm{x}) + b \right).
        \]
\end{itemize}
The second and third bullets expand \(\bm{w}\) by its representation
\(\bm{w} = \sum_n \alpha_n y_n \bm{z}_n\) from \autoref{sec:dual-messages},
so \(\bm{w}\) itself never has to be stored.

\begin{moral}
    \textbf{Kernel trick}: plug in an efficient kernel function to avoid
    dependence on \(\tilde{d}\).
\end{moral}

\subsection{Kernel SVM with QP}

\begin{algorithm}[H]
    \caption{Kernel Hard-Margin SVM}\label{alg:kernel-hard-svm}
    \(q_{n,m} = y_n y_m K(\bm{x}_n, \bm{x}_m)\); \(\bm{p} = -\bm{1}_N\);
        \((\mathrm{A}, \bm{c})\) for equ./bound constraints\;
    \(\bm{\alpha} \gets \mathrm{QP}(\mathrm{Q}_{\mathrm{D}}, \bm{p},
        \mathrm{A}, \bm{c})\)\;
    \(b \gets \left( y_s - \sum_{\text{SV indices } n} \alpha_n y_n
        K(\bm{x}_n, \bm{x}_s) \right)\) with SV \((\bm{x}_s, y_s)\)\;
    \Return SVs and their \(\alpha_n\) as well as \(b\) such that for new
        \(\bm{x}\),
        \(g_{\mathrm{SVM}}(\bm{x}) = \sign \left( \sum_{\text{SV indices } n}
        \alpha_n y_n K(\bm{x}_n, \bm{x}) + b \right)\)\;
\end{algorithm}

\begin{itemize}
    \item Line 1: time complexity \(O(N^2) \cdot\) (kernel evaluation).
    \item Line 2: QP with \(N\) variables and \(N + 1\) constraints.
    \item Lines 3 \& 4: time complexity \(O(\#\text{SV}) \cdot\) (kernel
        evaluation).
\end{itemize}

\(\tilde{d}\) appears nowhere: it is hidden inside the cost of one kernel
evaluation, which for \(\bm{\Phi}_2\) is \(O(d)\). The returned model is the
list of SVs with their \(\alpha_n\), not a weight vector.

\begin{moral}
    \textbf{Kernel SVM}: use a computational shortcut to avoid
    \(\tilde{d}\) \& predict with SVs only.
\end{moral}

\begin{exercise}[Fun Time]
    Consider two examples \(\bm{x}\) and \(\bm{x}'\) such that
    \(\bm{x}^\top \bm{x}' = 10\). What is \(K_{\bm{\Phi}_2}(\bm{x},
    \bm{x}')\)?
    \begin{enumerate}[(1)]
        \item \(1\);
        \item \(11\);
        \item \(111\);
        \item \(1111\).
    \end{enumerate}
\end{exercise}

\begin{proof}[Reference Answer]
    (3). Using the derivation in the previous subsections,
    \(K_{\bm{\Phi}_2}(\bm{x}, \bm{x}') = 1 + \bm{x}^\top \bm{x}' +
    (\bm{x}^\top \bm{x}')^2 = 1 + 10 + 100\).
\end{proof}

\section{Polynomial Kernel}\label{sec:poly-kernel}

\source{slides p.7--11}

\subsection{General Poly-2 Kernel}

Rescaling the coordinates of \(\bm{\Phi}_2\) changes its inner product, and
some rescalings give a tidier kernel:
\[
    \begin{aligned}
        \bm{\Phi}_2(\bm{x}) &= (1, x_1, \dots , x_d, x_1^2, \dots , x_d^2)
            &&\iff K_{\bm{\Phi}_2}(\bm{x}, \bm{x}') = 1 + \bm{x}^\top \bm{x}'
            + (\bm{x}^\top \bm{x}')^2\\
        \bm{\Phi}_2(\bm{x}) &= (1, \sqrt{2} x_1, \dots , \sqrt{2} x_d, x_1^2,
            \dots , x_d^2) &&\iff K_2(\bm{x}, \bm{x}') = 1 + 2 \bm{x}^\top
            \bm{x}' + (\bm{x}^\top \bm{x}')^2\\
        \bm{\Phi}_2(\bm{x}) &= (1, \sqrt{2\gamma} x_1, \dots , \sqrt{2\gamma}
            x_d, \gamma x_1^2, \dots , \gamma x_d^2) &&\iff K_2(\bm{x},
            \bm{x}') = 1 + 2\gamma \bm{x}^\top \bm{x}' + \gamma^2
            (\bm{x}^\top \bm{x}')^2
    \end{aligned}
\]
(the `\(\dots\)' after \(x_1^2\) still run through all the products \(x_i
x_j\)). The last line is a perfect square.

\begin{moral}
    \[
        K_2(\bm{x}, \bm{x}') = (1 + \gamma \bm{x}^\top \bm{x}')^2 \quad
        \text{with } \gamma > 0
    \]
\end{moral}

\begin{itemize}
    \item \(K_2\): somewhat `easier' to calculate than \(K_{\bm{\Phi}_2}\)
        --- one inner product, one addition, one square.
    \item \(\bm{\Phi}_2\) and the rescaled \(\bm{\Phi}_2\): equivalent
        \textbf{power}, different \textbf{inner product} \(\implies\)
        different \textbf{geometry}.
\end{itemize}
Every rescaled transform spans the same quadratic curves, so the hypothesis
sets coincide. But the SVM picks the hyperplane with the largest margin, and
the margin is measured with the inner product of \(\mathcal{Z}\); stretching
some coordinates of \(\mathcal{Z}\) changes which hyperplane is fattest.

\begin{moral}
    \(K_2\) is commonly used.
\end{moral}

\subsection{Poly-2 Kernels in Action}

\begin{center}
    \begin{tabular}{@{}c@{\hspace{1.2em}}c@{\hspace{1.2em}}c@{}}
        \ksvmpanel[3.6]{\ksvmPolyA}{\ksvmPolyData}{\ksvmPolyASV} &
        \ksvmpanel[3.6]{\ksvmPolyB}{\ksvmPolyData}{\ksvmPolyBSV} &
        \ksvmpanel[3.6]{\ksvmPolyC}{\ksvmPolyData}{\ksvmPolyCSV}\\[4pt]
        \((1 + 0.001 \bm{x}^\top \bm{x}')^2\) &
        \(1 + \bm{x}^\top \bm{x}' + (\bm{x}^\top \bm{x}')^2\) &
        \((1 + 1000 \bm{x}^\top \bm{x}')^2\)
    \end{tabular}
\end{center}

Each picture is a hard-margin SVM solved on the same data, with the
support vectors boxed. The pictures are mine: I measured the 40 examples off
the deck's figure and solved the dual with each kernel. The deck does not
say where its origin is, so the boundaries and SV sets differ in detail from
the deck's. I put the origin at the centre of the hyperbola in the deck's
third picture and doubled the scale.
\begin{itemize}
    \item \(g_{\mathrm{SVM}}\) \textbf{different}, \fbox{SVs}
        \textbf{different} --- `hard' to say which is better before learning.
        The three runs have \(4\), \(3\) and \(4\) SVs, and only the third
        grows a second branch at the bottom.
    \item Change of \textbf{kernel} \(\iff\) change of \textbf{margin
        definition}. The three margins, \(1/\norm{\bm{w}}\) in each kernel's
        own \(\mathcal{Z}\), come out as about \(0.0002\), \(0.15\) and
        \(120\): they are not even on a common scale.
\end{itemize}

\begin{moral}
    Need to select \(K\), just like selecting \(\bm{\Phi}\).
\end{moral}

\subsection{General Polynomial Kernel}

\[
    \begin{aligned}
        K_2(\bm{x}, \bm{x}') &= (\zeta + \gamma \bm{x}^\top \bm{x}')^2
            \text{ with } \gamma > 0, \zeta \ge 0\\
        K_3(\bm{x}, \bm{x}') &= (\zeta + \gamma \bm{x}^\top \bm{x}')^3
            \text{ with } \gamma > 0, \zeta \ge 0\\
        &\ \ \vdots\\
        K_Q(\bm{x}, \bm{x}') &= (\zeta + \gamma \bm{x}^\top \bm{x}')^Q
            \text{ with } \gamma > 0, \zeta \ge 0
    \end{aligned}
\]

\begin{center}
    \begin{minipage}[c]{0.58\textwidth}
        \begin{itemize}
            \item embeds \(\bm{\Phi}_Q\) specially with parameters \((\gamma,
                \zeta)\);
            \item allows computing large-margin \textbf{polynomial}
                classification \textbf{without dependence on
                \(\tilde{d}\)}.
        \end{itemize}
    \end{minipage}%
    \hfill
    \begin{minipage}[c]{0.38\textwidth}
        \centering
        \ksvmpanel[3.4]{\ksvmPolyTen}{\ksvmGaussData}{\ksvmPolyTenSV}\\[2pt]
        {\small 10-th order polynomial with margin 0.1}
    \end{minipage}
\end{center}

Expanding \((\zeta + \gamma \bm{x}^\top \bm{x}')^Q\) by the binomial
theorem gives every monomial up to degree \(Q\), each with its own positive
weight. So \(K_Q\) is the inner product of a rescaled \(\bm{\Phi}_Q\), whose
\(\tilde{d} = O(d^Q)\) never has to be paid. In two dimensions a 10th-order
transform has 66 coordinates, and the kernel still costs one inner product
in \(\mathbb{R}^2\) and one power.

The picture on the right is my own run of \((1 + \bm{x}^\top \bm{x}')^{10}\)
on the data of the Gaussian pictures below, in a frame of my choosing. Its
margin comes out as about \(0.30\) rather than the deck's \(0.1\), which is
another reminder that the margin depends on the frame and the kernel.

\begin{moral}
    SVM \(+\) Polynomial Kernel: \textbf{Polynomial SVM}.
\end{moral}

\subsection{Special Case: Linear Kernel}

\[
    \begin{aligned}
        K_1(\bm{x}, \bm{x}') &= (0 + 1 \cdot \bm{x}^\top \bm{x}')^1\\
        &\ \ \vdots\\
        K_Q(\bm{x}, \bm{x}') &= (\zeta + \gamma \bm{x}^\top \bm{x}')^Q
            \text{ with } \gamma > 0, \zeta \ge 0
    \end{aligned}
\]

\begin{center}
    \begin{minipage}[c]{0.62\textwidth}
        \begin{itemize}
            \item \(K_1\): just the \textbf{usual inner product}, called the
                \vocab{linear kernel}.
            \item `Even easier': can be solved (often in primal form)
                \textbf{efficiently}.
        \end{itemize}
    \end{minipage}%
    \hfill
    \begin{minipage}[c]{0.34\textwidth}
        \centering
        \svmtoypanel[band]{1}{-1}{-1}
    \end{minipage}
\end{center}

With \(\bm{\Phi} = \) identity, \(\tilde{d} = d\) and the kernel SVM is the
linear SVM of \autoref{sec:svm} --- the picture is its toy example again.
When \(d\) is small the primal QP with \(d + 1\) variables is the cheaper of
the two forms.

\begin{moral}
    Linear first, remember? :-)
\end{moral}

\begin{exercise}[Fun Time]
    Consider the general 2-nd polynomial kernel \(K_2(\bm{x}, \bm{x}') =
    (\zeta + \gamma \bm{x}^\top \bm{x}')^2\). Which of the following
    transforms can be used to derive this kernel?
    \begin{enumerate}[(1)]
        \item \(\bm{\Phi}(\bm{x}) = (1, \sqrt{2\gamma} x_1, \dots ,
            \sqrt{2\gamma} x_d, \gamma x_1^2, \dots , \gamma x_d^2)\);
        \item \(\bm{\Phi}(\bm{x}) = (\zeta, \sqrt{2\gamma} x_1, \dots ,
            \sqrt{2\gamma} x_d, x_1^2, \dots , x_d^2)\);
        \item \(\bm{\Phi}(\bm{x}) = (\zeta, \sqrt{2\gamma\zeta} x_1, \dots ,
            \sqrt{2\gamma\zeta} x_d, x_1^2, \dots , x_d^2)\);
        \item \(\bm{\Phi}(\bm{x}) = (\zeta, \sqrt{2\gamma\zeta} x_1, \dots ,
            \sqrt{2\gamma\zeta} x_d, \gamma x_1^2, \dots , \gamma x_d^2)\).
    \end{enumerate}
\end{exercise}

\begin{proof}[Reference Answer]
    (4). Since \((\zeta + \gamma \bm{x}^\top \bm{x}')^2 = \zeta^2 + 2\gamma
    \zeta \bm{x}^\top \bm{x}' + \gamma^2 (\bm{x}^\top \bm{x}')^2\), we need
    \(\zeta^2\) from the 0-th order terms, \(2\gamma\zeta \bm{x}^\top
    \bm{x}'\) from the 1-st order terms, and \(\gamma^2 (\bm{x}^\top
    \bm{x}')^2\) from the 2-nd order terms.
\end{proof}

\section{Gaussian Kernel}\label{sec:gaussian-kernel}

\source{slides p.12--16}

\subsection{Kernel of Infinite Dimensional Transform}

A polynomial kernel is cheap however large \(Q\) is. Pushing further: an
infinite dimensional \(\bm{\Phi}(\bm{x})\)? Yes, if \(K(\bm{x}, \bm{x}')\)
is efficiently computable! Take \(\bm{x} = (x)\), one dimension:
\[
    \begin{aligned}
        K(x, x') &= \exp(-(x - x')^2)\\
        &= \exp(-(x)^2) \exp(-(x')^2) \exp(2xx')\\
        &\overset{\text{Taylor}}{=} \exp(-(x)^2) \exp(-(x')^2) \left(
            \sum_{i=0}^\infty \frac{(2xx')^i}{i!} \right)\\
        &= \sum_{i=0}^\infty \left( \exp(-(x)^2) \exp(-(x')^2)
            \sqrt{\frac{2^i}{i!}} \sqrt{\frac{2^i}{i!}} (x)^i (x')^i
            \right)\\
        &= \bm{\Phi}(x)^\top \bm{\Phi}(x')
    \end{aligned}
\]
with the infinite dimensional
\[
    \bm{\Phi}(x) = \exp(-x^2) \cdot \left( 1, \sqrt{\frac{2}{1!}} x,
    \sqrt{\frac{2^2}{2!}} x^2, \dots \right).
\]
The middle step splits \((x - x')^2 = x^2 - 2xx' + x'^2\); the Taylor series
of \(\exp(2xx')\) is a polynomial kernel of every degree at once, and the
two outer factors attach \(\exp(-x^2)\) to each side.

\begin{definition}[Gaussian kernel]
    More generally, the \vocab{Gaussian kernel} is
    \[
        K(\bm{x}, \bm{x}') = \exp \left( -\gamma \norm{\bm{x} - \bm{x}'}^2
        \right) \quad \text{with } \gamma > 0.
    \]
\end{definition}

\subsection{Hypothesis of Gaussian SVM}

With the Gaussian kernel \(K(\bm{x}, \bm{x}') = \exp \left( -\gamma
\norm{\bm{x} - \bm{x}'}^2 \right)\),
\[
    \begin{aligned}
        g_{\mathrm{SVM}}(\bm{x}) &= \sign \left( \sum_{\text{SV}} \alpha_n
            y_n K(\bm{x}_n, \bm{x}) + b \right)\\
        &= \sign \left( \sum_{\text{SV}} \alpha_n y_n \exp \left( -\gamma
            \norm{\bm{x} - \bm{x}_n}^2 \right) + b \right).
    \end{aligned}
\]
\begin{itemize}
    \item A linear combination of Gaussians \textbf{centered at the SVs}
        \(\bm{x}_n\).
    \item Also called the \vocab{Radial Basis Function (RBF)} kernel.
\end{itemize}
Each SV contributes a bump of height \(\alpha_n\), pointing up for \(y_n =
+1\) and down for \(y_n = -1\), with width set by \(\gamma\). The bump
depends on \(\bm{x}\) only through its distance to \(\bm{x}_n\), which is
why the kernel is called radial.

\begin{moral}
    \textbf{Gaussian SVM}: find \(\alpha_n\) to combine Gaussians centered at
    \(\bm{x}_n\) \& achieve large margin in an infinite-dim.\ space.
\end{moral}

\subsection{Support Vector Mechanism}

\begin{center}
    \solidrules
    \begin{tabular}{c|c}
        \toprule
        & large-margin hyperplanes \(+\) higher-order transforms with
            kernel trick\\
        \midrule
        \# & \textbf{not many}\\
        boundary & \textbf{sophisticated}\\
        \bottomrule
    \end{tabular}
\end{center}

\begin{itemize}
    \item transformed vector \(\bm{z} = \bm{\Phi}(\bm{x}) \implies\)
        efficient kernel \(K(\bm{x}, \bm{x}')\);
    \item store optimal \(\bm{w} \implies\) store a few SVs and
        \(\alpha_n\).
\end{itemize}

This is the table that closed \autoref{sec:why-large-margin}, now with the
`numerous feature transform' made affordable. Neither the transform nor the
weight vector is ever stored, which is what makes an infinite-dimensional
\(\mathcal{Z}\) possible at all.

\begin{moral}
    New possibility by Gaussian SVM: \textbf{infinite-dimensional linear
    classification}, with generalization `guarded by' large-margin :-)
\end{moral}

\subsection{Gaussian SVM in Action}

\begin{center}
    \begin{tabular}{@{}c@{\hspace{1.2em}}c@{\hspace{1.2em}}c@{}}
        \ksvmpanel[3.6]{\ksvmGaussA}{\ksvmGaussData}{\ksvmGaussASV} &
        \ksvmpanel[3.6]{\ksvmGaussB}{\ksvmGaussData}{\ksvmGaussBSV} &
        \ksvmpanel[3.6]{\ksvmGaussC}{\ksvmGaussData}{\ksvmGaussCSV}\\[4pt]
        \(\exp\bigl(-1 \norm{\bm{x} - \bm{x}'}^2\bigr)\) &
        \(\exp\bigl(-10 \norm{\bm{x} - \bm{x}'}^2\bigr)\) &
        \(\exp\bigl(-100 \norm{\bm{x} - \bm{x}'}^2\bigr)\)
    \end{tabular}
\end{center}

These are again my runs, on 40 examples measured off the deck's figure and
placed in \([-1, 1]^2\). As \(\gamma\) grows, each Gaussian narrows and the
boundary bends around individual examples. At \(\gamma = 100\) it has broken
into islands around the crosses, and all \(40\) examples are SVs (\(9\) and
\(25\) for the other two).
\begin{itemize}
    \item Large \(\gamma \implies\) sharp Gaussians \(\implies\)
        `overfit'?
    \item \textbf{Warning: SVM can still overfit :-(}
\end{itemize}
The large-margin guard of \autoref{sec:why-large-margin} restricts the
dichotomies for a fixed \(\mathcal{Z}\), but \(\gamma\) changes
\(\mathcal{Z}\). The right-hand picture separates the data perfectly and
still looks like memorization.

\begin{moral}
    Gaussian SVM: need \textbf{careful selection of \(\gamma\)}.
\end{moral}

\begin{exercise}[Fun Time]
    Consider the Gaussian kernel \(K(\bm{x}, \bm{x}') = \exp\bigl(-\gamma
    \norm{\bm{x} - \bm{x}'}^2\bigr)\). What function does the kernel converge to
    if \(\gamma \to \infty\)?
    \begin{enumerate}[(1)]
        \item \(K_{\lim}(\bm{x}, \bm{x}') = 0\);
        \item \(K_{\lim}(\bm{x}, \bm{x}') = \llbracket \bm{x} = \bm{x}'
            \rrbracket\);
        \item \(K_{\lim}(\bm{x}, \bm{x}') = \llbracket \bm{x} \neq \bm{x}'
            \rrbracket\);
        \item \(K_{\lim}(\bm{x}, \bm{x}') = 1\).
    \end{enumerate}
\end{exercise}

\begin{proof}[Reference Answer]
    (2). If \(\bm{x} = \bm{x}'\), \(K(\bm{x}, \bm{x}') = 1\) regardless of
    \(\gamma\). If \(\bm{x} \neq \bm{x}'\), \(K(\bm{x}, \bm{x}') = 0\) when
    \(\gamma \to \infty\). Thus \(K_{\lim}\) is an impulse function, which is
    an extreme case of how the Gaussian gets sharper when \(\gamma \to
    \infty\). With it, \(g_{\mathrm{SVM}}\) returns the sign of \(b\)
    everywhere except exactly at the training inputs.
\end{proof}

\section{Comparison of Kernels}\label{sec:kernel-comparison}

\source{slides p.17--21}

\subsection{Linear Kernel: Cons and Pros}

\begin{center}
    \begin{minipage}[c]{0.4\textwidth}
        \centering
        \svmtoypanel[band]{1}{-1}{-1}
    \end{minipage}%
    \hfill
    \begin{minipage}[c]{0.5\textwidth}
        \centering
        \(K(\bm{x}, \bm{x}') = \bm{x}^\top \bm{x}'\)
    \end{minipage}
\end{center}

\begin{center}
    \begin{minipage}[t]{0.47\textwidth}
        \raggedright\textbf{Cons}
        \begin{itemize}
            \item restricted --- \textbf{not always separable?!}
        \end{itemize}
        \centering
        \begin{tikzpicture}[scale=1.1]
            \circdataX
            \xplotframe
        \end{tikzpicture}
    \end{minipage}%
    \hfill
    \begin{minipage}[t]{0.47\textwidth}
        \raggedright\textbf{Pros}
        \begin{itemize}
            \item safe --- \textbf{linear first, remember? :-)}
            \item fast --- with a \textbf{special QP solver} in primal
            \item very explainable --- \textbf{\(\bm{w}\) and SVs} say
                something
        \end{itemize}
    \end{minipage}
\end{center}

The `not separable' example is the circle data set of
\autoref{sec:nonlin-transform}: no line splits the circles inside from the
crosses outside, and a hard-margin linear SVM has no solution there.
Explainable means that each \(w_i\) weighs one original feature, as in the
credit example of \autoref{sec:perceptron-set}.

\begin{moral}
    Linear kernel: an important \textbf{basic} tool.
\end{moral}

\subsection{Polynomial Kernel: Cons and Pros}

\begin{center}
    \begin{minipage}[c]{0.4\textwidth}
        \centering
        \ksvmpanel[2.8]{\ksvmPolyA}{\ksvmPolyData}{\ksvmPolyASV}
    \end{minipage}%
    \hfill
    \begin{minipage}[c]{0.5\textwidth}
        \centering
        \(K(\bm{x}, \bm{x}') = (\zeta + \gamma \bm{x}^\top \bm{x}')^Q\)
    \end{minipage}
\end{center}

\begin{center}
    \begin{minipage}[t]{0.47\textwidth}
        \raggedright\textbf{Cons}
        \begin{itemize}
            \item \textbf{numerical difficulty} for large \(Q\):
                \begin{itemize}
                    \item \(\abs{\zeta + \gamma \bm{x}^\top \bm{x}'} < 1\):
                        \(K \to 0\);
                    \item \(\abs{\zeta + \gamma \bm{x}^\top \bm{x}'} > 1\):
                        \(K \to\) big;
                \end{itemize}
            \item three parameters \((\gamma, \zeta, Q)\) --- \textbf{more
                difficult to select}.
        \end{itemize}
    \end{minipage}%
    \hfill
    \begin{minipage}[t]{0.47\textwidth}
        \raggedright\textbf{Pros}
        \begin{itemize}
            \item \textbf{less restricted} than linear;
            \item strong physical control --- `knows' \textbf{degree \(Q\)}.
        \end{itemize}
    \end{minipage}
\end{center}

The numerical problem is the \(Q\)-th power: entries of
\(\mathrm{Q}_{\mathrm{D}}\) spread over many orders of magnitude, and the QP
solver loses precision. The \((1 + 1000 \bm{x}^\top \bm{x}')^2\) run above
already shows it: its margin of about \(120\) against \(0.0002\) for
\(\gamma = 0.001\) reflects kernel values of wildly different sizes.

\begin{moral}
    Polynomial kernel: perhaps \textbf{small-\(Q\) only} --- sometimes
    efficiently done by \textbf{linear on \(\bm{\Phi}_Q(\bm{x})\)}.
\end{moral}

\subsection{Gaussian Kernel: Cons and Pros}

\begin{center}
    \begin{minipage}[c]{0.4\textwidth}
        \centering
        \ksvmpanel[2.8]{\ksvmGaussA}{\ksvmGaussData}{\ksvmGaussASV}
    \end{minipage}%
    \hfill
    \begin{minipage}[c]{0.5\textwidth}
        \centering
        \(K(\bm{x}, \bm{x}') = \exp\bigl(-\gamma \norm{\bm{x} - \bm{x}'}^2\bigr)\)
    \end{minipage}
\end{center}

\begin{center}
    \begin{minipage}[t]{0.47\textwidth}
        \raggedright\textbf{Cons}
        \begin{itemize}
            \item \textbf{mysterious} --- no \(\bm{w}\);
            \item \textbf{slower} than linear;
            \item \textbf{too powerful?!}
        \end{itemize}
        \centering
        \ksvmpanel[2.4]{\ksvmGaussC}{\ksvmGaussData}{\ksvmGaussCSV}
    \end{minipage}%
    \hfill
    \begin{minipage}[t]{0.47\textwidth}
        \raggedright\textbf{Pros}
        \begin{itemize}
            \item \textbf{more powerful} than linear/poly.;
            \item bounded --- \textbf{less numerical difficulty} than poly.;
            \item one parameter only --- \textbf{easier to select} than poly.
        \end{itemize}
    \end{minipage}
\end{center}

`No \(\bm{w}\)' because \(\bm{w}\) lives in an infinite-dimensional
\(\mathcal{Z}\) and can only be held as its SV representation. `Bounded'
because \(0 < K \le 1\) for every pair. `Too powerful' is the
\(\gamma = 100\) picture again.

\begin{moral}
    Gaussian kernel: \textbf{one of the most popular} but shall be
    \textbf{used with care}.
\end{moral}

\subsection{Other Valid Kernels}

\begin{itemize}
    \item A kernel represents a \textbf{special} similarity:
        \(\bm{\Phi}(\bm{x})^\top \bm{\Phi}(\bm{x}')\).
    \item Any similarity \(\implies\) valid kernel? \textbf{Not really.}
    \item Necessary \& sufficient conditions for a valid kernel:
        \vocab{Mercer's condition}
        \begin{itemize}
            \item symmetric;
            \item let \(k_{ij} = K(\bm{x}_i, \bm{x}_j)\); the matrix
                \(\mathrm{K}\)
                \[
                    \begin{aligned}
                        \mathrm{K} &= \begin{bmatrix}
                            \bm{\Phi}(\bm{x}_1)^\top \bm{\Phi}(\bm{x}_1) &
                            \bm{\Phi}(\bm{x}_1)^\top \bm{\Phi}(\bm{x}_2) &
                            \dots & \bm{\Phi}(\bm{x}_1)^\top
                            \bm{\Phi}(\bm{x}_N)\\
                            \bm{\Phi}(\bm{x}_2)^\top \bm{\Phi}(\bm{x}_1) &
                            \bm{\Phi}(\bm{x}_2)^\top \bm{\Phi}(\bm{x}_2) &
                            \dots & \bm{\Phi}(\bm{x}_2)^\top
                            \bm{\Phi}(\bm{x}_N)\\
                            \dots & \dots & \dots & \dots\\
                            \bm{\Phi}(\bm{x}_N)^\top \bm{\Phi}(\bm{x}_1) &
                            \bm{\Phi}(\bm{x}_N)^\top \bm{\Phi}(\bm{x}_2) &
                            \dots & \bm{\Phi}(\bm{x}_N)^\top
                            \bm{\Phi}(\bm{x}_N)
                        \end{bmatrix}\\
                        &= \begin{bmatrix} \bm{z}_1 & \bm{z}_2 & \dots &
                            \bm{z}_N \end{bmatrix}^\top \begin{bmatrix}
                            \bm{z}_1 & \bm{z}_2 & \dots & \bm{z}_N
                            \end{bmatrix}\\
                        &= \mathrm{Z} \mathrm{Z}^\top
                    \end{aligned}
                \]
                must \textbf{always} be \textbf{positive semi-definite}.
        \end{itemize}
\end{itemize}

Necessity is one line: for any \(\bm{v} \in \mathbb{R}^N\), \(\bm{v}^\top
\mathrm{Z} \mathrm{Z}^\top \bm{v} = \norm{\mathrm{Z}^\top \bm{v}}^2 \ge 0\).
Here \(\mathrm{Z}\) has the \(\bm{z}_n^\top\) as rows, as in the foundations.
Sufficiency --- that a symmetric \(K\) whose matrices are always PSD
\emph{is} an inner product in some \(\mathcal{Z}\) --- is Mercer's theorem,
which we claim. The condition must hold for every choice of inputs
\(\bm{x}_1, \dots , \bm{x}_N\), not just the training set.

\begin{moral}
    Define your own kernel: possible, \textbf{but hard}.
\end{moral}

\begin{exercise}[Fun Time]
    Which of the following is \textbf{not} a valid kernel? (Hint: consider
    two 1-dimensional vectors \(\bm{x}_1 = (1)\) and \(\bm{x}_2 = (-1)\) and
    check Mercer's condition.)
    \begin{enumerate}[(1)]
        \item \(K(\bm{x}, \bm{x}') = (-1 + \bm{x}^\top \bm{x}')^2\);
        \item \(K(\bm{x}, \bm{x}') = (0 + \bm{x}^\top \bm{x}')^2\);
        \item \(K(\bm{x}, \bm{x}') = (1 + \bm{x}^\top \bm{x}')^2\);
        \item \(K(\bm{x}, \bm{x}') = (-1 - \bm{x}^\top \bm{x}')^2\).
    \end{enumerate}
\end{exercise}

\begin{proof}[Reference Answer]
    (1). The kernels in (2) and (3) are just polynomial kernels. The kernel
    in (4) is equivalent to the kernel in (3), since \((-1 - s)^2 = (1 +
    s)^2\). For (1), the matrix formed from the kernel and the two examples
    is
    \[
        \mathrm{K} = \begin{bmatrix} (-1 + 1)^2 & (-1 - 1)^2 \\ (-1 - 1)^2 &
        (-1 + 1)^2 \end{bmatrix} = \begin{bmatrix} 0 & 4 \\ 4 & 0
        \end{bmatrix},
    \]
    with eigenvalues \(\pm 4\), so it is not positive semi-definite. Thus the
    underlying kernel is not a valid one.
\end{proof}

\hrulebar

\subsection*{Summary}

\begin{description}
    \item[Kernel Trick] kernel as a shortcut of transform \(+\) inner
        product.
    \item[Polynomial Kernel] embeds a specially-scaled polynomial transform.
    \item[Gaussian Kernel] embeds an infinite dimensional transform.
    \item[Comparison of Kernels] linear for efficiency or Gaussian for
        power.
\end{description}

The kernel removed every trace of \(\tilde{d}\), and with it the last brake
on the transform: the Gaussian SVM at \(\gamma = 100\) separates anything
and overfits. The hard margin demands that every example be classified
correctly, which is exactly what a noisy data set should not get. Next:
avoiding overfitting in the Gaussian (and other kernels) by letting the
margin be violated.
